\chapter{What is claimed, what is not, and what comes next}
\markboth{Limits and next steps}{Limits and next steps}
\section{What is shown}
Memory, audit, reach and fate for ten soils in a one-dimensional matrix-flow model. The design modules reproduce published worked examples. Parsers and shared functions agree with independent readers, and defects were removed. Daisy 7.1.14 runs end to end from SWAFL, including a 5.85 ha field from a satellite image to a schedule. Data, code, manual and application are open.

\section{What is not claimed}
\begin{itemize}[leftmargin=*]
\item Field validation: measured drainage, lysimeters and sensor-scheduled plots are still needed. The memory result gives a testable prediction: sand should drain about 85 per cent of the released storage and silty clay loam about 2 per cent.
\item Run-by-run repetition of the published experiments.
\item The Windows build and the online data services on a connected machine.
\item The non-equilibrium criterion $\Pi_\tau$ is proposed, not tested.
\item Several model terms need attention: the baseline offset, sprinkler range coefficients, the wind term and the conductivity floor.
\end{itemize}
The constant $\kappa$ describes the crop, the climate and the sensor, and all E2 runs used one cereal rotation, one weather series and a sensor at 20 cm. For another crop or climate, $\kappa$ is determined again from a trigger experiment on a few soils.

\section{What is covered and what is not yet}
Part~I derives the equations of the manuscripts; Part~II works the mathematics of the SWAFL layers with computed examples: Engine~1 and Engine~2, irrigation control, the field balance and the soil--plant--atmosphere resistances, soil structure and tillage, mulch and evaporation, feedback and the hidden questions, phenotyping, surface, drip and sprinkler design, drainage design, the crop library and crop water requirement, weather and field data (evapotranspiration, satellite indices, runoff and soil loss), the Dashboard, and the governing equations of the Daisy workbench (Richards column, roots, heat, solutes, nitrogen, crop, calibration). Not yet worked in this edition: the reports, workflows and validity chapters of the manual; the 107 reference topics, 69 experiment factors and 43 best combinations of the Daisy workbench (tables of Daisy parameters and runs); the screenshot atlas; the full numerical analysis of the Richards solver (only one worked column is shown); nitrogen transformations beyond the forms in Chapter~\ref{ch:daisy}; and a collection of laboratory problems on the open ensemble. The book has not reached the length first planned; it holds what has been worked and checked.
